\section{期间送达}
\begin{enumerate}
  \item 考期间的计算。
  \item 一般规则。
        \begin{enumerate}
          \item 月从当月某日到下月某日为 1 个月，没有就到下个月最后一日，不一定要凑够 30 天。
          \item 半个月就是 15 天，但 1 个月未必是 30 天。
          \item 刑期提前一天放出去。
          \item 顺延到节假日后的第一个工作日。羁押实打实算。
          \item 法定期间不包括在途时间。
          \item 交邮应当以邮件上的\zd{邮戳为证}。
        \end{enumerate}
  \item 特殊规则。
        \begin{enumerate}
          \item 恢复计算。\zd{前提是过期}。\al{9 月 21 日一审宣判，当庭送达判决书，9 月 30 日上诉书交看守所，10 月 8 日上班时才寄出，10 月 10 日寄到法院}。
                \begin{enumerate}
                  \item 当事人去申请；正当事由 \hz 不可抗拒的事由；及时申请，障碍消除后的 5 天内。
                  \item 最终还是法院说了算，耽误了几天给你补几天。
                \end{enumerate}
          \item 重新计算。归零。\zd{机关改变了、程序要重来}。仅包括公检法的办案期限。
          \item 不计算。\zd{身份不明阅卷}、诉讼\zd{中}止，\zd{精神鉴定监委送}。\zd{回避属于延期审理，要计入审限}。
        \end{enumerate}
  \item 送达同民诉。
  \item 公安机关重新计算侦查羁押期限的，应当报批准逮捕的检察院备案。
\end{enumerate}
